% Tracce uscite agli Esami di Stato (Sezione A, Settore Informazione) sugli argomenti del capitolo.
% Incluso da main.tex subito dopo il riassunto.

\section*{Tracce Uscite} % Usiamo * per non numerare la sezione nell'indice
\addcontentsline{toc}{section}{Tracce Uscite: Programmazione e Algoritmi}

\subsection*{Prima sessione 2016 -- Informatica}
\textbf{Traccia}: concetto di classe, di oggetto e differenze nella programmazione ad oggetti; esempio in pseudocodice o in un linguaggio a scelta di una classe \texttt{Rational} (numeratore e denominatore interi) con funzione \texttt{mult} (moltiplicazione per un intero) che ritorna un nuovo oggetto \texttt{Rational}.

\textbf{Dove trovare la risposta}: sezione \emph{Classi e Oggetti}, con la classe \texttt{Rational} completa e le motivazioni (immutabilità, invarianti).

\subsection*{Prima sessione 2018 -- Informatica (prima prova)}
\textbf{Traccia}: implementare \texttt{fact(x)}, una versione iterativa e una ricorsiva, discutendo le differenze tra i due approcci.

\textbf{Dove trovare la risposta}: sezione \emph{Ricorsione e Iterazione}.

\subsection*{Prima sessione 2018 -- Informatica (seconda prova)}
\textbf{Traccia}: algoritmo di corrispondenza tra due immagini $W\times H$, con ricerca dello spostamento $d \in [0, D]$; complessità in tempo e memoria.

\textbf{Dove trovare la risposta}: sezione \emph{Tema Svolto: Corrispondenze tra Due Immagini}.

\subsection*{Seconda sessione 2018 -- Informatica}
\textbf{Traccia}: modello ISO/OSI con ogni livello e confronto con TCP/IP.

\textbf{Dove trovare la risposta}: capitolo \emph{Reti di Calcolatori}.

\subsection*{Prima sessione 2024 -- Informatica \#3}
\textbf{Traccia}: ricerca informata e non informata, con efficienza, memoria e ottimalità.

\textbf{Dove trovare la risposta}: capitolo \emph{Intelligenza Artificiale e Machine Learning}; per BFS, DFS e Dijkstra vedi la sezione \emph{Strutture Dati e Grafi}.

\subsection*{Prima sessione 2025 (Bologna) -- Ing. Informatica E05}
\textbf{Traccia}: processo di chiamata a funzione; ruolo dello stack nell'allocazione dei dati; distinzione semantica tra parametri formali e attuali; conseguenze delle diverse modalità di passaggio dei parametri; esempi di codice in un linguaggio a piacere o pseudocodice.

\textbf{Dove trovare la risposta}: sezione \emph{Chiamate di Funzione, Stack e Passaggio dei Parametri}.

\subsection*{Prima sessione 2026 -- A1-INF-4}
\textbf{Traccia}: polimorfismo nella programmazione orientata agli oggetti: principi, finalità e vantaggi (modularità, estensibilità, riusabilità); modalità di realizzazione e ruolo nell'interazione tra interfacce e implementazioni; limiti e problematiche; esempio applicativo nell'Ingegneria dell'Informazione con motivazione delle scelte progettuali.

\textbf{Dove trovare la risposta}: sezione \emph{Polimorfismo}, con la scaletta di risposta e l'esempio dei sensori.

\subsection*{Bologna (sessione non indicata nel file) -- Ing. Informatica E04}
\textbf{Traccia}: concetti di processo e thread in un sistema operativo moderno; differenze; modelli di multithreading (many-to-one, one-to-one, many-to-many); problematiche di concorrenza come le race condition.

\textbf{Dove trovare la risposta}: capitolo \emph{Sistemi Operativi}, sezione \emph{Processi e Thread}; per le race condition e la sincronizzazione vedi anche la sezione \emph{Concorrenza e Integrità} del capitolo \emph{Progettazione di Sistemi e API}.
